\begingroup
\setlength{\parindent}{0pt}
\setlength{\parskip}{0.58\baselineskip}
\setlength{\abovedisplayskip}{9pt plus 2pt minus 2pt}
\setlength{\belowdisplayskip}{9pt plus 2pt minus 2pt}
\setlength{\jot}{7pt}

\section*{How the Proof Works}

We are asking whether a viscous fluid that starts smooth can ever stop being
smooth.  This proof answers no.  The difficulty is that viscosity is not the
only part of the equation changing the fluid.  The fluid also carries and
deforms its own velocity:
\[
\underbrace{\partial_tu}_{\text{local acceleration}}
+\underbrace{(u\!\cdot\nabla)u}_{\text{nonlinear transport}}
=-\underbrace{\nabla p}_{\text{pressure}}
+\underbrace{\nu\Delta u}_{\text{viscous diffusion}},
\qquad \nabla\!\cdot u=0.
\]
The nonlinear term can sharpen the flow while viscosity smooths it.  Pressure
keeps the motion incompressible and couples every local change to the
surrounding field.  The regularity problem is decided by whether nonlinear
sharpening can outrun viscous smoothing before some finite time \(T_*\).

The place to see that danger is a vortex being stretched.  Picture the
surrounding flow pulling it lengthwise.  Incompressibility conserves the volume
of a material tube, so its cross-section contracts as its length increases.
The core narrows.  If the velocity change across that core does not fall just
as quickly, then
\[
|\nabla u|\sim\frac{|\Delta U|}{\ell}
\]
grows as the radius \(\ell\) shrinks.  The fluid turns more sharply across
less distance: rotation strengthens and the core steepens.  This is nonlinear
sharpening in the fluid itself.  Viscosity resists it by diffusing velocity
across the narrowing core.

Now the mathematics can name the contest.  Write
\[
\omega:=\nabla\times u,
\qquad
S:=\frac12\bigl(\nabla u+(\nabla u)^{\mathsf T}\bigr).
\]
Taking the curl of Navier--Stokes gives
\[
(\partial_t+u\!\cdot\nabla)\omega
=S\omega+\nu\Delta\omega.
\]
When the vorticity points along an extensional direction of \(S\), the
stretching term \(S\omega\) strengthens it.  Viscosity acts against that
sharpening through \(\nu\Delta\omega\).  After integration over space, the
contest is visible in one line:
\[
\frac12\frac{d}{dt}\|\omega(t)\|_2^2
=\int_{\mathbb R^3}\omega\!\cdot S\omega\,dx
-\nu\|\nabla\omega(t)\|_2^2.
\]
The first term is nonlinear production.  The second is viscous removal.

Nothing here requires the kinetic energy itself to become infinite.  In fact,
\[
\frac12\frac{d}{dt}\|u(t)\|_2^2
+\nu\|\nabla u(t)\|_2^2=0.
\]
The total energy decreases.  The unresolved possibility is concentration:
finite energy can be carried by thinner regions while the velocity gradients
inside those regions grow.  A vortex can be losing energy and still be
sharpening near its core.

So we shrink the picture and ask whether the balance changes.  For
\(\lambda>0\), set
\[
u_\lambda(x,t)=\lambda u(\lambda x,\lambda^2t),
\qquad
p_\lambda(x,t)=\lambda^2p(\lambda x,\lambda^2t).
\]
Every part of the equation picks up one common factor:
\[
\begin{aligned}
\partial_tu_\lambda
  &=\lambda^3(\partial_tu)(\lambda x,\lambda^2t),
&
(u_\lambda\!\cdot\nabla)u_\lambda
  &=\lambda^3((u\!\cdot\nabla)u)(\lambda x,\lambda^2t),\\
\nabla p_\lambda
  &=\lambda^3(\nabla p)(\lambda x,\lambda^2t),
&
\nu\Delta u_\lambda
  &=\lambda^3\nu(\Delta u)(\lambda x,\lambda^2t).
\end{aligned}
\]
Viscosity does not acquire an extra power of \(\lambda\) that makes it win
automatically in a smaller vortex.  Nonlinear transport, pressure, local
acceleration, and diffusion remain matched by scaling.

The corresponding Sobolev calculation tells us where that match lives.  Since
\[
\|u_\lambda(\cdot,t)\|_{\dot H^s}
=\lambda^{s-1/2}\|u(\cdot,\lambda^2t)\|_{\dot H^s},
\]
concentration lowers the norm below \(s=\tfrac12\), raises it above
\(s=\tfrac12\), and leaves it unchanged at
\[
s=\frac12.
\]
This is the critical height.  It is not a claim that nonlinear production and
viscous dissipation are numerically equal at each moment.  It says that scale
alone gives neither one an advantage there.

Now keep following the vortex as its core narrows.  If its radius is
\(\ell\), a second spatial derivative across that core has size
\(\ell^{-2}\) relative to the velocity change it sees.  The viscous term acts
on the corresponding interval
\[
|\nu\Delta u|
\sim\frac{\nu}{\ell^2}|u|,
\qquad
\tau_\ell\sim\frac{\ell^2}{\nu}.
\]
If the core kept narrowing through
\(\ell_j=\rho^j\ell_0\) with \(0<\rho<1\), those intervals would satisfy
\[
\sum_{j=0}^{\infty}\frac{\ell_j^2}{\nu}
=\frac{\ell_0^2}{\nu}\sum_{j=0}^{\infty}\rho^{2j}
=\frac{\ell_0^2}{\nu(1-\rho^2)}<\infty.
\]
That is the Zeno picture inside the vortex: each thinner core asks the equation
to respond faster, and the resulting intervals can fit below one finite
terminal time.  It does not assert that a solution follows this cascade.  It
shows the failure pattern the proof has to prevent.

A finite maximal endpoint makes the next step exact.  If \(T_*<\infty\), the
endpoint continuation criterion requires
\[
\sup_{0<t<T_*}\|u(t)\|_{L^3}=\infty.
\]
Here velocity blow-up means escape of the global critical \(L^3\) norm; it
does not mean that every particle speed must diverge.  For every finite temporal
order \(j\),
\[
\partial_t^ju(t)-\partial_t^ju(0)
=\int_0^t\partial_s^{j+1}u(s)\,ds.
\]
Taking the \(L^3\) norm gives
\[
\|\partial_t^ju(t)\|_3
\leq \|\partial_t^ju(0)\|_3
  +\int_0^t\|\partial_s^{j+1}u(s)\|_3\,ds.
\]
The smooth datum supplies every finite initial row.  If one later row stayed
uniformly bounded on the finite interval, repeated integration would bound
every row below it, including \(u\), contradicting the required \(L^3\)
escape.  Thus, for every fixed \(j\),
\[
\sup_{0<t<T_*}\|\partial_t^ju(t)\|_3=\infty.
\]

That is the velocity--acceleration--jerk insight in plain language.  In order
for a singularity to form, velocity has to blow up.  In order for
that to happen, its acceleration would also have to blow up, which means its
jerk has to blow up, and so on, and so on:
\[
u,\qquad \partial_tu,\qquad \partial_t^2u,\qquad\ldots
\]
This is the turn in the argument.  Acceleration is already written into
Navier--Stokes.  Once we substitute the equation for it and differentiate, the
jerk appears.  Differentiate again and the next response appears.  The first
three rungs are
\[
\begin{aligned}
\partial_tu
&=\nu\Delta u-(u\!\cdot\nabla)u-\nabla p,\\[4pt]
\partial_t^2u
&=\nu\Delta\partial_tu
  -(\partial_tu\!\cdot\nabla)u
  -(u\!\cdot\nabla)\partial_tu
  -\nabla\partial_tp,\\[4pt]
\partial_t^3u
&=\nu\Delta\partial_t^2u
  -(\partial_t^2u\!\cdot\nabla)u
  -2(\partial_tu\!\cdot\nabla)\partial_tu\\
&\hspace{2.6em}
  -(u\!\cdot\nabla)\partial_t^2u
  -\nabla\partial_t^2p.
\end{aligned}
\]
The complete pattern is
\[
\partial_t^{r+1}u
=\nu\Delta\partial_t^ru
-\sum_{a=0}^{r}\binom ra
  (\partial_t^au\!\cdot\nabla)\partial_t^{r-a}u
-\nabla\partial_t^rp.
\]
Only now do we compress the notation:
\[
V_r:=\partial_t^ru,
\qquad
Q_r:=\partial_t^rp.
\]
Pressure has to travel with every response.  The compact tower is
\[
\begin{aligned}
Q_r
&=R_iR_j\sum_{a=0}^{r}\binom ra V_{a,i}V_{r-a,j},\\
V_{r+1}
&=\nu\Delta V_r
  -\sum_{a=0}^{r}\binom ra(V_a\!\cdot\nabla)V_{r-a}
  -\nabla Q_r.
\end{aligned}
\]

Now the connection to spatial smoothness is visible.  A hypothetical
singularity runs forward through rougher and faster temporal responses.  The
proof reads the viscous relation backward.  Once a finite response and its
transport--pressure terms are controlled, the equation isolates
\(\Delta V_r\) and returns two spatial derivatives to the preceding response.
Each finite climb in time opens a finite climb in Sobolev height when traced
back to the original velocity.  As the temporal order rises, the Sobolev
height available to \(u\) rises with it.

Now return to the critical height that scaling exposed.  Let
\[
\Lambda:=(-\Delta)^{1/2},
\qquad
E_s(t):=\frac12\|\Lambda^su(t)\|_2^2,
\qquad
H(t):=E_{1/2}(t).
\]
Pairing Navier--Stokes at height \(s\) gives
\[
E_s'=\mathcal P_s-2\nu E_{s+1},
\]
where \(\mathcal P_s\) is the complete signed nonlinear production at that
height.  Repeated differentiation at \(s=\tfrac12\) now gives
\[
\begin{aligned}
H'
&=\mathcal P_{1/2}-2\nu E_{3/2},\\
H''
&=\partial_t\mathcal P_{1/2}
  -2\nu\mathcal P_{3/2}
  +(2\nu)^2E_{5/2},\\
H'''
&=\partial_t^2\mathcal P_{1/2}
  -2\nu\partial_t\mathcal P_{3/2}
  +(2\nu)^2\mathcal P_{5/2}
  -(2\nu)^3E_{7/2}.
\end{aligned}
\]
At every finite order,
\[
H^{(n)}
=(-2\nu)^nE_{n+1/2}
+\sum_{j=0}^{n-1}(-2\nu)^j
 \partial_t^{\,n-1-j}\mathcal P_{j+1/2}.
\]
The \(n\)-th temporal derivative has exposed the spatial energy at height
\(n+\tfrac12\).  The calculation lifts directly to every temporal response:
\[
E_{r,s}(t):=\frac12\|\Lambda^sV_r(t)\|_2^2,
\qquad
\partial_tE_{r,s}
=\mathcal P_{r,s}-2\nu E_{r,s+1}.
\]
If every finite height is controlled, then any requested classical derivative
follows by choosing \(n\) high enough:
\[
H^{n+1/2}(\mathbb R^3)\hookrightarrow C_b^k(\mathbb R^3)
\qquad\text{whenever }n>k+1,
\]
and the full spatial intersection gives
\[
\bigcap_{m<\infty}H^m(\mathbb R^3)
=H^\infty(\mathbb R^3)
\hookrightarrow C^\infty(\mathbb R^3).
\]
This is the realization behind the tower.  It is not yet the bound.  The
differentiated identities show which spatial coordinates have to be paid for;
they do not make those coordinates finite by rearrangement.

The payment is the whole-terminal rectangle theorem.  Fix any finite temporal
depth \(N\) and spatial height \(M\).  For each \(0\le r\le N\), it supplies
\[
V_r\in L^2(0,T_*;H^{M+1}),
\qquad
V_{r+1}\in L^2(0,T_*;H^{M-1}),
\]
together with the pressure coordinates carried by the recurrence.  The
Hilbert-pair trace step then gives
\[
V_r\in C([0,T_*];H^M).
\]
This is one finite rectangle.  Its constant may depend on \(N\) and \(M\).
The proof never asks one constant to control infinitely many derivatives at
once.  It proves the required statement for every finite rectangle.

Those rectangles are not unrelated solutions.  Every coordinate belongs to
one velocity history generated by \(u_0\).  When two rectangles overlap, their
shared coordinates agree.  Compatibility lets us assemble the entire family:
\[
u_0
\longrightarrow
\{\mathcal R_{N,M}[u_0;T_*]\}_{N,M<\infty}
\longrightarrow
\mathcal I[u_0]
:=\bigcap_{r,m<\infty}H^r(0,T_*;H^m(\mathbb R^3)),
\qquad
u\in\mathcal I[u_0].
\]
No infinite-order estimate has been smuggled into this step.  Each coordinate
came from a finite theorem; the intersection records that all of them belong
to one compatible fluid.

Now take the zeroth temporal row.  For every finite \(m\), the rectangle pair
and the equation place
\[
u\in C([0,T_*];H^m_\sigma).
\]
All of these traces meet at one terminal value,
\[
u_*:=u(T_*)
\in\bigcap_{m<\infty}H^m_\sigma
=H^\infty_\sigma
\hookrightarrow C^\infty.
\]
At \(T_*\), the equation converts the completed spatial information back into
time information.  Starting from \(V_0^*=u_*\), the pressure formula gives
\(Q_0^*\), and the recurrence gives \(V_1^*\).  Repeating this at every finite
order produces the pressure-complete terminal jet.  Smooth local existence
restarts from \(u_*\); uniqueness joins that continuation to the original
history.  This extends the solution past \(T_*\), contradicting finite
maximality.  Hence
\[
T_*=\infty.
\]

Once the intersection is present, the familiar estimates are read from it
instead of proved one by one.  Given a finite response order \(r\), a spatial
derivative order \(k\), and \(2\le p\le\infty\), choose
\[
m>k+\frac32-\frac3p.
\]
That finite coordinate gives
\[
\mathcal I[u_0]
\longrightarrow H_t^1H_x^m
\longrightarrow C_tH_x^m
\longrightarrow L_t^qW_x^{k,p}
\qquad(1\le q\le\infty).
\]
Choosing \(m=1,2,3\) in turn recovers, among many others,
\[
L_t^\infty L_x^6,
\qquad
L_t^\infty L_x^\infty,
\qquad
L_t^\infty W_x^{1,\infty}.
\]
These are not separate achievements.  Each is one finite readout of
\(\mathcal I[u_0]\).

The whole motion is
\[
\boxed{
u_0
\longrightarrow \{(V_r,Q_r)\}_{r\ge0}
\longrightarrow \{\mathcal R_{N,M}[u_0]\}_{N,M<\infty}
\longrightarrow \mathcal I[u_0]
\longrightarrow u_*\in H^\infty
\longrightarrow \text{restart}
\longrightarrow T_*=\infty.
}
\]

\endgroup
